\chapter*{Résumé}
\addcontentsline{toc}{chapter}{Résumé}

L'accès du citoyen marocain à l'information juridique reste difficile : les textes sont
publics, mais rédigés dans une langue dense et parfois disponibles uniquement en arabe. Les
assistants conversationnels généralistes, eux, connaissent mal le droit marocain, peuvent
inventer des références inexistantes, et font transiter les questions des citoyens hors du
territoire national.

Ce stage a consisté à concevoir, réaliser et évaluer un assistant de recherche juridique
\textbf{souverain} fondé sur l'architecture RAG (\textit{Retrieval-Augmented Generation}),
appliqué au Code du travail marocain (Loi n\textsuperscript{o}~65-99). Le système répond à une
question posée en langage naturel en s'appuyant \emph{uniquement} sur les articles réellement
retrouvés dans le corpus, cite systématiquement leurs numéros, et s'abstient explicitement
lorsque l'information demandée n'y figure pas. L'ensemble de la chaîne --- corpus, embeddings,
recherche et génération --- s'exécute en local, sans aucun appel à une API externe.

Le corpus a été construit à la main à partir du PDF officiel : 589 articles extraits, nettoyés
et enrichis de leurs métadonnées de hiérarchie, avec une régénération vérifiée identique octet
pour octet depuis la source. La recherche combine une méthode sémantique (embeddings BGE-M3,
base vectorielle Chroma) et une méthode lexicale (BM25), fusionnées par \textit{Reciprocal Rank
Fusion}. La génération s'appuie sur un modèle Mistral~7B exécuté localement via Ollama, encadré
par deux garde-fous : un seuil d'abstention appliqué \emph{avant} l'appel au modèle, et une
vérification que chaque article cité figurait bien dans le contexte fourni.

Le système a été évalué sur un jeu de référence de 37 questions. La recherche atteint un
Recall@5 de 0{,}97 et un MRR de 0{,}89 en configuration dense. Côté réponses, 97~\% citent au
moins un article, 94~\% des citations sont vérifiées comme réellement fournies au modèle, et
l'abstention est correcte sur 100~\% des questions hors périmètre. Le temps de réponse médian
est de 3{,}5~secondes, réduit à 0{,}14~seconde lorsque le garde-fou d'abstention se déclenche.
Les limites de mesure --- notamment l'impossibilité de vérifier automatiquement qu'un article
cité \emph{dit réellement} ce que la réponse lui attribue --- sont documentées explicitement.

\noindent\rule[2pt]{\textwidth}{0.5pt}

{\textbf{Mots clés :}}
RAG, souveraineté numérique, traitement automatique du langage naturel, recherche sémantique,
grands modèles de langage, droit du travail marocain, ancrage documentaire, abstention.
\\
\noindent\rule[2pt]{\textwidth}{0.5pt}
